\section{US SBL Worked Example: SEC Rule 15c3-3 with Cash Collateral}
\label{sec:appI}
\label{app:us-sbl-example}

This appendix is a regulatory correctness witness for the US securities-lending regime: cash collateral, rebate mechanics, and the rehypothecation cap of SEC Rule~15c3-3~\cite{sec15c33}, with FINRA SLATE reporting mapped at each step. It is the counterpart to the European title-transfer example (App.~\ref{sec:appH}). Title transfer is not pledge; the two regimes exercise the same Ledger primitives under different law, and both witnesses are retained.

\subsection{Participants and Setup}

\begin{center}\small
\begin{tabular}{l l p{7.5cm}}
\toprule
\textbf{Entity} & \textbf{Role} & \textbf{Description} \\
\midrule
FIDELITY & Lender (beneficial owner) & US mutual fund. Holds 1{,}000{,}000 shares of NVIDIA (NVDA, ISIN US67066G1040). Agent lender: State~Street. \\
MORGAN\_STANLEY & Borrower & US broker-dealer. Borrows NVDA to cover a client short. Subject to Rule~15c3-3. \\
STATE\_STREET & Agent lender & Acts for FIDELITY; negotiates and manages the loan. \\
\bottomrule
\end{tabular}
\end{center}

\noindent\textbf{Loan terms.} 1{,}000{,}000 NVDA, open loan, cash collateral (USD), margin 102\%. Cash rebate: Fed Funds Effective Rate minus 25~bps. Take Fed Funds $= 5.25\%$, so rebate rate $= 5.00\%$. The lender reinvests the cash collateral and earns the spread between reinvestment return and rebate paid. As in App.~\ref{sec:appH}, intermediate quantities are computed in exact integer minor units; the fractional cash figures shown below (rebates, spreads) are display roundings, with round-half-to-even (IEEE~754) applied only at the settlement-instruction projection boundary.

\medskip
\noindent\textbf{Legal regime} \texttt{US\_15C3\_3}.
\begin{itemize}[nosep]
\item Cash collateral transfers from Morgan~Stanley to Fidelity/State~Street, who reinvest it.
\item Rehypothecation of customer securities is capped at 140\% of the customer's debit balance (Rule~15c3-3(b)(3)).
\item The poster's PnL includes $\mathrm{coll\_post}$: $V_e^{\mathrm{15c3}} = \sum_u (\vec{w}_e(u)[\mathrm{own}] + \vec{w}_e(u)[\mathrm{coll\_post}]) \cdot P_t(u)$.
\end{itemize}

\medskip
\noindent\textbf{Locate precondition (Reg~SHO Rule~203(b)(1))~\cite{regsho}.} Before the client's short sale executes, Morgan~Stanley satisfies the locate obligation, either by holding the borrow or by a locate confirmation against inventory. The locate is pre-trade validation: no moves occur, no coordinate changes, and conservation is not engaged until the borrow executes.

\subsection{Step 1: Loan Initiation (Day 0)}
\label{sec:us-step1}

NVDA price \$875.00. Loan value and required collateral:
\[
\mathrm{LV} = 1{,}000{,}000 \times 875 = \$875{,}000{,}000, \qquad
\mathrm{RCV} = 875{,}000{,}000 \times 102\% = \$892{,}500{,}000.
\]
Cash is always owned; the collateral leg writes $\mathrm{own}$ on the poster and $\mathrm{coll\_recv}$ on the receiver, the earmark that must be returned at close.

\noindent Each $\tau$ below is a single Transaction carrying moves, with $\mathtt{txIntroduce}=\textsc{Nothing}$ and $\mathtt{txAppend}=\textsc{Nothing}$: neither a registration nor an amendment. Its move list is the conserved flow, so conservation holds by construction from the signed edge-sum; the trailing comment names the economic role, not a type field on the core Transaction.

\begin{lstlisting}
tau_1 = Transaction(day 0):  -- securities-loan settlement
  -- Securities leg (unit: NVDA)
  Move 1: FIDELITY        onloan    += 1,000,000    [onloan]
  Move 2: MORGAN_STANLEY  borr      += 1,000,000    [borr]
  -- Cash collateral leg (unit: USD)
  Move 3: MORGAN_STANLEY  own       -= 892,500,000  [own]
  Move 4: FIDELITY        coll_recv += 892,500,000  [coll_recv]
\end{lstlisting}

\noindent\textbf{Units.} NVDA (equity), USD (cash). Position vector
\[(\mathrm{own},\,\mathrm{onloan},\,\mathrm{borr},\,\mathrm{coll\_post},\,\mathrm{coll\_recv},\,\mathrm{coll\_rehyp}),\]
abbreviated own, onl, borr, c\_p, c\_r, c\_rh.

\begin{center}\footnotesize\setlength{\tabcolsep}{4pt}
\begin{tabular}{l l r r r r r r l}
\toprule
\textbf{Entity} & \textbf{Unit} & own & onl & borr & c\_p & c\_r & c\_rh & avail \\
\midrule
FIDELITY & NVDA & 1{,}000{,}000 & 1{,}000{,}000 & 0 & 0 & 0 & 0 & 0 \\
FIDELITY & USD & 0 & 0 & 0 & 0 & 892{,}500{,}000 & 0 & 0 \\
\midrule
MORGAN\_STANLEY & NVDA & 0 & 0 & 1{,}000{,}000 & 0 & 0 & 0 & 1{,}000{,}000 \\
MORGAN\_STANLEY & USD & $-892{,}500{,}000$ & 0 & 0 & 0 & 0 & 0 & $-892{,}500{,}000$ \\
\bottomrule
\end{tabular}
\end{center}

\noindent\textbf{P18 check.} Fidelity $\vec{w}(\text{NVDA})[\mathrm{own}] = 1{,}000{,}000$, unchanged. \checkmark\quad
\textbf{SLATE.} New Loan Event: FINRA Loan ID, Client Loan ID, execution timestamp, ISIN/CUSIP, quantity, collateral type (Cash), collateral amount, rebate rate, term type (Open), borrower/lender LEI.

\subsection{Step 2: Rebate Mechanics and Cash Reinvestment}
\label{sec:us-step2}

Fidelity reinvests the \$892{,}500{,}000 in approved money-market instruments, earning the money-market yield (take $= 5.25\%$) and paying Morgan~Stanley the rebate ($5.00\%$). The reinvestment transforms the cash earmark into MMF units (see \S\ref{sec:sbl-reinvestment}):

\begin{lstlisting}
tau_reinvest = Transaction(day 0):  -- cash reinvestment
  -- Deploys the cash into MMF units; the USD coll_recv tables below
  -- carry the unchanged return obligation, with the MMF overlay separate.
  Move 1: FIDELITY  coll_recv -= 892,500,000  [coll_recv, USD]
  Move 2: FIDELITY  own       += 892,500       [own, MMF_UNIT]
  Move 3: MMF_FUND  own       -= 892,500       [own, MMF_UNIT]
  Move 4: MMF_FUND  own       += 892,500,000   [own, USD]
\end{lstlisting}

\noindent The reinvestment deploys the cash, but the \$892.5M return obligation persists in the loan unit state (\texttt{collateral\_amount}): the collateral must be returned in full at close. The position-vector tables that follow therefore carry $\mathrm{coll\_recv}$ (USD) at the collateral \emph{obligation} amount throughout---the cash-equivalent Fidelity owes back---while the 892{,}500 MMF units are held in parallel, in a separate unit, and tracked separately. The margin-call (Step~4) and release (Steps~6--7) moves accordingly act on the obligation amount shown. Fidelity now bears reinvestment risk; conservation holds in each unit independently.

\medskip
\noindent\textbf{Daily economics.}
\begin{align*}
\text{Reinvestment income} &= 892{,}500{,}000 \times 5.25\% / 360 = \$130{,}104.17 \\
\text{Rebate payable} &= 892{,}500{,}000 \times 5.00\% / 360 = \$123{,}958.33 \\
\text{Lending spread} &= 130{,}104.17 - 123{,}958.33 = \$6{,}145.83
\end{align*}

\noindent Billing here is monthly. Rebate at Day~30: $892{,}500{,}000 \times 5.00\% \times 30/360 = \$3{,}718{,}750$.

\begin{lstlisting}
tau_rebate = Transaction(day 30):  -- rebate (fee settlement)
  Move 1: FIDELITY        own -= 3,718,750   [own, USD]
  Move 2: MORGAN_STANLEY  own += 3,718,750   [own, USD]
\end{lstlisting}

\noindent Reinvestment over the month $= \$3{,}903{,}906$; rebate paid $= \$3{,}718{,}750$; net $= \$185{,}156$, the lending spread.

\subsection{Step 3: Rehypothecation Constraints under Rule 15c3-3}
\label{sec:us-step3}

Rule~15c3-3(b)(3) permits Morgan~Stanley to use customer securities only up to 140\% of the customer's aggregate debit balance. This contrasts with the title-transfer regime of App.~\ref{sec:appH}:

\begin{center}\small
\begin{tabular}{l p{5cm} p{5cm}}
\toprule
& \textbf{European (GMSLA~2010 TT)} & \textbf{US (Rule 15c3-3)} \\
\midrule
Rehypothecation & Receiver rehypothecates collateral freely & Borrower faces a regulatory cap \\
Legal basis & Title transfer: receiver \emph{owns} collateral & Pledge/control: customer retains ownership interest \\
Cap & None (SFTR Art.~15 documentation only) & 140\% of customer debit balance \\
\bottomrule
\end{tabular}
\end{center}

\noindent\textbf{Pre-condition enforcement (P19).} The SBL contract validates a rehypothecation request before any move commits; the \texttt{legal\_regime} field selects the path:

\begin{lstlisting}
validate_rehypothecation(entity, unit, quantity, regime):
  if regime == US_15C3_3:
    max_rehyp     = customer_debit_balance(entity) * 1.40
    current_rehyp = sum(entity.coll_rehyp[u] for u in units)
    if current_rehyp + quantity * price(unit) > max_rehyp:
      REJECT("15c3-3 rehypothecation cap exceeded")
  elif regime == TITLE_TRANSFER:
    if not sftr_art15_documented(entity): REJECT("SFTR Art 15 missing")
    PERMIT()
\end{lstlisting}

\noindent A request exceeding the cap is rejected, never committed: P19 is a pre-condition, so the illegal state is unreachable rather than reached and rolled back.

\subsection{Step 4: Daily Mark-to-Market (Day 5)}
\label{sec:us-step4}

NVDA rises to \$910.00.
\[
\mathrm{RCV}_5 = 1{,}000{,}000 \times 910 \times 102\% = \$928{,}200{,}000,
\]
\[
\Delta C = 928{,}200{,}000 - 892{,}500{,}000 = \$35{,}700{,}000.
\]

\begin{lstlisting}
tau_2 = Transaction(day 5):  -- margin call
  Move 1: MORGAN_STANLEY  own       -= 35,700,000   [own, USD]
  Move 2: FIDELITY        coll_recv += 35,700,000   [coll_recv, USD]
\end{lstlisting}

\begin{center}\small
\begin{tabular}{l l r r r r r r}
\toprule
\textbf{Entity} & \textbf{Unit} & own & onl & borr & c\_p & c\_r & c\_rh \\
\midrule
FIDELITY & NVDA & 1{,}000{,}000 & 1{,}000{,}000 & 0 & 0 & 0 & 0 \\
FIDELITY & USD & 0 & 0 & 0 & 0 & 928{,}200{,}000 & 0 \\
\midrule
MORGAN\_STANLEY & NVDA & 0 & 0 & 1{,}000{,}000 & 0 & 0 & 0 \\
MORGAN\_STANLEY & USD & $-928{,}200{,}000$ & 0 & 0 & 0 & 0 & 0 \\
\bottomrule
\end{tabular}
\end{center}

\noindent\textbf{P13 check.} Collateral held $= \$928{,}200{,}000 =$ required. \checkmark\quad
\textbf{SLATE.} Modification: collateral amount and value updated.

\subsection{Step 5: FINRA SLATE Reporting}
\label{sec:us-step5}

Rule~10c-1a~\cite{sec10c1a} and the FINRA Rule~6500 Series~\cite{finra6500} require every covered securities loan to be reported to SLATE. The lifecycle maps as follows.

\begin{center}\small
\begin{tabular}{l l p{6cm}}
\toprule
\textbf{Lifecycle Event} & \textbf{SLATE Event} & \textbf{Key Data Elements} \\
\midrule
Loan initiation (Day~0) & New Loan & FINRA Loan ID, Client Loan ID, timestamp, ISIN US67066G1040, quantity 1{,}000{,}000, collateral type Cash, amount \$892.5M, rebate 5.00\%, term Open, borrower/lender LEI \\
\midrule
Margin call (Day~5) & Modification & Collateral amount $\to$ \$928.2M; timestamp \\
\midrule
Partial return (Day~15) & Modification & Quantity $\to$ 500{,}000; collateral reduced \\
\midrule
Full return (Day~25) & Modification & Quantity $\to$ 0; loan closed \\
\bottomrule
\end{tabular}
\end{center}

\noindent\textbf{Deadlines.} Rule~10c-1a (adopted January 2024) and the FINRA Rule~6500 Series (effective 2025) set: loans effected on a business day from 12:00am~ET through 7:45pm~ET report same day before 8:00pm~ET; loans after 7:45pm~ET report next business day before 8:00pm~ET.

\noindent\textbf{Dual identifiers.} Each loan carries a FINRA Loan ID (assigned by SLATE) and a Client Loan ID (the firm's $\mathrm{loan\_id}$ in the SBL unit state).

\noindent\textbf{Unsettled Loan Flag (Field~\#44).} Set when the loan has not settled on the reporting day. The SBL settlement status (\S\ref{sec:sbl-state-machine}) maps directly: \texttt{PENDING}~$\to$~Unsettled, \texttt{ACTIVE}~$\to$~Settled.

\begin{remark}
The SBL settlement status is a materialised projection of the immutable event log---a read cache the log rebuilds---not an authoritative mutable store. Every status change is caused by a logged lifecycle event; replaying events to any point reconstructs the status that held then. The lifecycle authority is \S\ref{sec:lifecycle}.
\end{remark}

\subsection{Step 6: Partial Recall and Return (Day 15)}
\label{sec:us-step6}

Fidelity recalls 500{,}000 shares; NVDA \$920.00.
\[
\mathrm{RCV}_{15} = 500{,}000 \times 920 \times 102\% = \$469{,}200{,}000,
\]
\[
C_{\mathrm{release}} = 928{,}200{,}000 - 469{,}200{,}000 = \$459{,}000{,}000.
\]

\begin{lstlisting}
tau_3 = Transaction(day 15):  -- partial return
  -- Securities return (unit: NVDA)
  Move 1: MORGAN_STANLEY  borr      -= 500,000      [borr]
  Move 2: FIDELITY        onloan    -= 500,000      [onloan]
  -- Cash collateral release (unit: USD)
  Move 3: FIDELITY        coll_recv -= 459,000,000  [coll_recv]
  Move 4: MORGAN_STANLEY  own       += 459,000,000  [own]
\end{lstlisting}

\begin{center}\footnotesize\setlength{\tabcolsep}{4pt}
\begin{tabular}{l l r r r r r r l}
\toprule
\textbf{Entity} & \textbf{Unit} & own & onl & borr & c\_p & c\_r & c\_rh & avail \\
\midrule
FIDELITY & NVDA & 1{,}000{,}000 & 500{,}000 & 0 & 0 & 0 & 0 & 500{,}000 \\
FIDELITY & USD & 0 & 0 & 0 & 0 & 469{,}200{,}000 & 0 & 0 \\
\midrule
MORGAN\_STANLEY & NVDA & 0 & 0 & 500{,}000 & 0 & 0 & 0 & 500{,}000 \\
MORGAN\_STANLEY & USD & $-469{,}200{,}000$ & 0 & 0 & 0 & 0 & 0 & $-469{,}200{,}000$ \\
\bottomrule
\end{tabular}
\end{center}

\noindent\textbf{P15 check.} Loan quantity $1{,}000{,}000 \to 500{,}000$, monotonically non-increasing. \checkmark\quad
\textbf{P18 check.} Fidelity $\vec{w}(\text{NVDA})[\mathrm{own}] = 1{,}000{,}000$, unchanged. \checkmark

\noindent Pro-rated rebate on the returned portion: $459{,}000{,}000 \times 5.00\% \times 15/360 = \$956{,}250$.

\begin{lstlisting}
tau_3b = Transaction(day 15):  -- pro-rated rebate (fee settlement)
  Move 1: FIDELITY        own -= 956,250    [own, USD]
  Move 2: MORGAN_STANLEY  own += 956,250    [own, USD]
\end{lstlisting}

\subsection{Step 7: Full Return and Close-Out (Day 25)}
\label{sec:us-step7}

Remaining 500{,}000 NVDA returned; NVDA \$905.00.

\begin{lstlisting}
tau_4 = Transaction(day 25):  -- full return
  -- Securities return (unit: NVDA)
  Move 1: MORGAN_STANLEY  borr      -= 500,000       [borr]
  Move 2: FIDELITY        onloan    -= 500,000       [onloan]
  -- Cash collateral release (unit: USD)
  Move 3: FIDELITY        coll_recv -= 469,200,000   [coll_recv]
  Move 4: MORGAN_STANLEY  own       += 469,200,000   [own]
\end{lstlisting}

\noindent Final rebate: $469{,}200{,}000 \times 5.00\% \times 10/360 = \$651{,}666.67$.

\begin{lstlisting}
tau_4b = Transaction(day 25):  -- final rebate (fee settlement)
  Move 1: FIDELITY        own -= 651,666.67   [own, USD]
  Move 2: MORGAN_STANLEY  own += 651,666.67   [own, USD]
\end{lstlisting}

\noindent The loan unit is disposed: $\vec{w}_F(u_{\mathrm{loan}})[\mathrm{own}] = \vec{w}_M(u_{\mathrm{loan}})[\mathrm{own}] = 0$; SBL state RETURNED.

\begin{center}\small
\begin{tabular}{l l r r r r r r l}
\toprule
\textbf{Entity} & \textbf{Unit} & own & onl & borr & c\_p & c\_r & c\_rh & avail \\
\midrule
FIDELITY & NVDA & 1{,}000{,}000 & 0 & 0 & 0 & 0 & 0 & 1{,}000{,}000 \\
FIDELITY & USD & $-1{,}607{,}917$ & 0 & 0 & 0 & 0 & 0 & $-1{,}607{,}917$ \\
\midrule
MORGAN\_STANLEY & NVDA & 0 & 0 & 0 & 0 & 0 & 0 & 0 \\
MORGAN\_STANLEY & USD & $+1{,}607{,}917$ & 0 & 0 & 0 & 0 & 0 & $+1{,}607{,}917$ \\
\bottomrule
\end{tabular}
\end{center}

\noindent Fidelity's negative USD is the total rebate paid ($956{,}250 + 651{,}666.67 = \$1{,}607{,}917$). Reinvestment income (\S\ref{sec:us-step2}) offsets it but accrues in MMF units, a separate unit, not the USD position shown; Fidelity's net loan income is reinvestment return minus rebate paid.

\medskip
\noindent\textbf{Conservation.}
\begin{itemize}[nosep]
\item NVDA: $1{,}000{,}000_F + 0_M + \text{virtual}(-1{,}000{,}000) = 0$. \checkmark
\item USD: $(-1{,}607{,}917)_F + 1{,}607{,}917_M = 0$ (plus virtual offsets for the original collateral movements). \checkmark
\end{itemize}

\noindent\textbf{SLATE.} Final Modification, quantity~$=0$, closing the loan.

\subsection{Cross-Jurisdictional Comparison}
\label{sec:sbl-comparison}

The two worked examples differ only in three places: the \texttt{legal\_regime} field of the SBL unit state, the smart-contract pre-conditions (P19), and the PnL projection. The six-coordinate position vector, the Single-Coordinate Move Principle, and the conservation law are identical across regimes.

\begin{center}\small
\begin{longtable}{p{3.2cm} p{5.2cm} p{5.2cm}}
\toprule
\textbf{Aspect} & \textbf{European (GMSLA~2010 TT)} & \textbf{US (Rule 15c3-3)} \\
\midrule
\endfirsthead
\toprule
\textbf{Aspect} & \textbf{European (GMSLA~2010 TT)} & \textbf{US (Rule 15c3-3)} \\
\midrule
\endhead
\bottomrule
\endfoot
Collateral ownership & Title transfers to receiver, who \emph{owns} collateral. & Pledge: poster retains economic ownership; receiver holds under control. \\
\midrule
Collateral type & Non-cash (German Bunds); haircut/margin on collateral market value. & Cash (USD); 102\% of loan value posted. \\
\midrule
Rehypothecation & Unrestricted under title transfer; SFTR Art.~15 documentation required. & Capped at 140\% of customer debit balance (Rule 15c3-3(b)(3)); contract blocks excess. \\
\midrule
Fee structure & Lending fee (25~bps on loan value); borrower pays lender. & Rebate (Fed Funds $-$ 25~bps); lender pays borrower on cash, earns reinvestment spread. \\
\midrule
Collateral coordinate & $\mathrm{own}$ for poster, $\mathrm{coll\_recv}$ for receiver. & $\mathrm{own}$ for poster, $\mathrm{coll\_recv}$ for receiver. \\
\midrule
PnL formula & $V_e = \sum_u \vec{w}_e(u)[\mathrm{own}] \cdot P_t(u)$ & $V_e = \sum_u (\vec{w}_e(u)[\mathrm{own}] + \vec{w}_e(u)[\mathrm{coll\_post}]) \cdot P_t(u)$ \\
\midrule
Regulatory reporting & SFTR (NEWT, MODI, VALU, ETRM); dual-sided; UTI per ESMA waterfall. & FINRA SLATE (New Loan, Modification); single-sided (lender reports). \\
\midrule
Manufactured dividends & French withholding-tax treatment. & US tax-treaty provisions (IRC \S871(m)). \\
\midrule
Settlement discipline & CSDR: daily cash penalties for fails; auto-partial (IBP-125). & Reg~SHO Rule~204: close-out by T+3; FINRA Rule~4320. \\
\midrule
Governing agreement & GMSLA~2010 (title transfer); ISLA Best Practice Handbook. & Master Securities Loan Agreement; SEC/FINRA rules. \\
\midrule
Position vector & Identical six coordinates. & Identical six coordinates. \\
\midrule
Conservation law & $\sum_e \vec{w}_e(u)[\text{coord}] = 0$. & Identical. \\
\midrule
Smart contract & Same SBL contract; \texttt{legal\_regime} selects validation. & Same SBL contract; \texttt{legal\_regime} selects validation. \\
\bottomrule
\end{longtable}
\end{center}

\noindent The Ledger engine, conservation law, and move primitive are jurisdiction-neutral. Title transfer versus pledge, unrestricted versus capped rehypothecation, and lending fee versus rebate are encoded as SBL unit-state parameters and enforced by P19 pre-conditions; only the PnL projection and the pre-condition logic differ, as derived in \S\ref{sec:sbl-title-transfer}.
